\subsection*{4.3 Factor reconstruction for ABP}
We show how to get the absolute factor $g$ of $f$ from the lifted factor. This is called the jump step in Sudan's lecture notes~\cite{ref18}. The difference to the earlier presentations is that our lifted factor might not be monic.
Let $f=gh$, where $f$ has degree $d$, factor $g$ is irreducible and coprime to $h$, and $f,g,h$ are monic in $x$. In the previous section, we started with a factorization $f\equiv g_0h_0\pmod{I_0}$, where $g_0$ is irreducible and coprime to $h_0$. Moreover,$g\equiv g_0h'_0\pmod{I_0}$, for some $h'_0$ such that $h_0=hh'_0\pmod{I_0}$.
Then we apply Hensel lifting, say $t$-times, for some $t$ to be determined below. By Lemma 13, we get a factorization $f\equiv g_th_t\pmod{I_t}$ such that
\begin{equation}g\equiv g_th'_t\pmod{I_t},\tag{11}\end{equation}
for some $h'_t$ such that $h_t\equiv hh'_t\pmod{I_t}$.
Equation(11) gives us a relation between the known $g_t$ and the unknown $g$, via the unknown $h'_t$. We set up a linear system of equations to find a polynomial $\widetilde g\in K[x,y]$ that is monic in $x$ and has minimal degree in $x$ such that
\begin{equation}\widetilde g\equiv g_t\widetilde h\pmod{I_t},\tag{12}\end{equation}
for some polynomial $\widetilde h$. We give some more details to the linear system below, after the next lemma. Before, we show that $\widetilde g$ is indeed the factor we were looking for, for large enough $t$.
\begin{ntlemma}[Author Lemma 16]
$\widetilde g=g$, for $t\geq\log(4d^2+1)$.
\end{ntlemma}
\begin{proof}
Consider the resultant $r(y)=\Res_x(g,\widetilde g)$. We show that $r(y)=0$. Then it follows from Lemma 2 that $g$ and $\widetilde g$ share a common factor with positive $x$-degree. Since $g$ is irreducible it must be a divisor of $\widetilde g$. Since both of them are monic and have the same $x$-degree, we get equality $\widetilde g=g$, up to constant multiples.
To argue that $r(y)=0$, recall from Lemma 2 that the resultant can be written as $r(y)=ug+v\widetilde g$, for some polynomials $u$ and $v$. By(11) and(12), we have
\[ug+v\widetilde g\equiv g_t(uh'_t+v\widetilde h)\pmod{I_t}.\]
Consider $g_t$ and $w=uh'_t+v\widetilde h$ as polynomials in $y$ with coefficients in $x$. Suppose $g_t=c_0(x)+c_1(x)y+\cdots+c_{d'}(x)y^{d'}$. By the properties of Hensel lifting, we have $g_t\equiv g_0\pmod{I_0}$, and therefore $c_0(x)=g_0(x)$. Recall that $g_0$ is non-constant,$\deg(g_0)\geq1$.
Let $j\geq0$ be the least power of $y$ that appears in $w$ and let its coefficient be $w_j(x)$. Suppose for the sake of contradiction that $j<2^t$. Then the least power of $y$ in $g_tw$ is also $j$, and its coefficient is $g_0(x)w_j(x)$, which is a nonzero polynomial in $x$.
The monomials present in $g_0(x)w_j(x)y^j$ cannot be canceled by other monomials in $g_tw$ because they have larger $y$-degree. It follows that $g_tw\bmod I_t$ is not free of $x$. On the other hand,$r(y)\equiv g_tw\pmod{I_t}$ and $r(y)\in K[y]$ is a polynomial with no variable $x$. This is a contradiction.
We conclude that $j\geq2^t$, which means that $w\equiv0\pmod{I_t}$. Hence, we get $r(y)\equiv0\pmod{I_t}$. Recall that $\deg(r)\leq4d^2$. Now we choose $t\geq\log(4d^2+1)$. Then we can conclude that indeed $r(y)=0$.
\end{proof}
\paragraph{Details for setting up the linear system}
Equation(12) can be used to set up a homogeneous system of linear equations. For the degree bounds of the polynomials, let $d_x=\deg_x(g)$ and $d_y=\deg_y(g)$. We may assume that we know $d_x$ and $d_y$. Let $D_x=\deg_x(g_t)$ and $D_y=4d^2$, where $d=\deg(f)$. Let $D'_x=\deg_x(h_t)$. Recall from Lemma 13 that $\deg_x(\widetilde h)\leq D'_x$. Let
\begin{align*}g_t&=\sum_{i\leq D_x,j\leq D_y}c_{i,j}x^iy^j,\\\widetilde g&=x^{d_x}+\sum_{i<d_x,j\leq d_y}r_{i,j}x^iy^j,\\\widetilde h&=\sum_{i\leq D'_x,j\leq D_y}s_{i,j}x^iy^j,\end{align*}
where the coefficients $c_{i,j},r_{i,j},s_{i,j}$ are polynomials in the variables $z_1,\ldots,z_n$. To ensure that $\widetilde g$ is monic, we set the coefficient of $x^{d_x}$ in $\widetilde g$ to be1.
Note that we have an ABP that computes $g_t$. Hence, there are ABPs for computing the coefficients $c_{i,j}$ of $g_t$ by Lemma 4. The coefficients $r_{i,j},s_{i,j}$ of $\widetilde g$ and $\widetilde h$ we treat as unknowns. Equation(12) now becomes
\begin{equation}\sum_{i\leq d_x,j\leq d_y}r_{i,j}x^iy^j\equiv\left(\sum_{i\leq D_x,j\leq D_y}c_{i,j}x^iy^j\right)\left(\sum_{i\leq D'_x,j\leq D_y}s_{i,j}x^iy^j\right)\pmod{y^{2d^2+1}}.\tag{13}\end{equation}
Now we equate the coefficients of the monomials $x^ky^\ell$ on both sides in(13). Then we get $(D_x+D'_x+1)(D_y+1)$ homogeneous linear equations in $d_x(d_y+1)+(D'_x+1)(D_y+1)$ many unknowns $r_{i,j}$ and $s_{i,j}$. This system can be expressed in the form $M\mathbf v=0$, for a matrix $M$ and unknown vector $\mathbf v$.
By Lemma 9, an ABP can efficiently compute a solution vector $\mathbf v$ of polynomials from $\mathbb F[\mathbf z]$. Note that by(11), a nontrivial solution is guaranteed to exist. From $\mathbf v$, we get the coefficients $r_{i,j}$ of $\widetilde g$, and hence of the factor $g$.
\subsection*{4.4 Size Analysis}
We summarize the bound on the ABP-size of the factor computed. Given polynomial $p$ of degree $d_p$ and $\sz(p)=s$. We have seen that the pre-processing transformations yield a polynomial $f$ of degree $d_f\leq2d_p$ and $\sz(f)=\operatorname{poly}(s)$. Then we do $t=\log(2d_f^2+1)$ iterations of Hensel lifting. The initial polynomials $f_0,g_0,h_0$ have ABP-size bound by $2d_f$. Hence, the polynomials after the last iteration have ABP-size bounded by $2^t\operatorname{poly}(s)=\operatorname{poly}(s,d_p)=\operatorname{poly}(s)$.
From the lifted factor we construct the actual factor of $f$. This step involves solving a linear system. We argued that the resulting polynomial $g$ has ABP-size $\operatorname{poly}(s)$.
Finally, we reverse the transformations from the beginning and get a factor of $p$ that has an ABP of size $\operatorname{poly}(s)$. This finishes the proof of Theorem 10.
